% §8.3 The Entropy Floor for Approximate Sets
%
% Monograph register: the -log2(eps) bits/element floor for FNR-zero sets, with
% the pigeonhole proof. Correctness fixes: (1) the two-sided lower bound is OPEN
% (source rem:general_lower_bound) -- the chapter previously asserted a specific
% "exact two-sided bound" n(1-H(omega))log2(1/eps), which is not a valid bound;
% (2) the per-element cost grows ADDITIVELY (+log2 10 per decade), not by
% doubling. Cross-refs sec:one-sided-two-sided, def:one-sided, def:bloom-filter,
% sec:bloom-gap, ch:optimal-bernoulli-sets. DEFINED: sec:entropy-floor,
% thm:entropy-floor, eq:entropy-floor, eq:accepted-size, eq:pigeonhole.

\section{The Entropy Floor for Approximate Sets}
\label{sec:entropy-floor}

The bit-budget comparison of \Cref{sec:one-sided-two-sided} placed the Bloom
filter at roughly $9.6$ bits per item at $\fprate = 0.01$ against a floor near
$6.64$. This section turns that contrast into a theorem: any approximate set
with $\fnrate = 0$ and false-positive rate at most $\fprate$, in any
representation whatever, must store at least $\log_2(1/\fprate)$ bits per
element on average. The Bloom filter pays a constant factor above this floor;
chapter~14 meets it.

\subsection*{Setup}

Fix a universe $\Universe$ with $|\Universe| = N$ and a target $A \subseteq
\Universe$ with $|A| = n$. An \emph{approximate-set representation} is a pair
$(R, \mathrm{query})$ with $R \in \{0,1\}^L$ a bit string and
$\mathrm{query}(R, \cdot) : \Universe \to \{\top, \bot\}$ a deterministic
membership algorithm; $L$ is the cost to bound. Following \Cref{def:one-sided}
we restrict to the FNR-zero family containing the Bloom filter:
$\mathrm{query}(R, x) = \top$ for every $x \in A$ (no inserted element is
missed), and at most a $\fprate$ fraction of the $N - n$ non-members is
accepted. The \emph{accepted set} $A_R = \{x : \mathrm{query}(R, x) = \top\}$
then satisfies $A \subseteq A_R$ and $|A_R \setminus A| \le \fprate(N - n)$, so
\begin{equation}
\label{eq:accepted-size}
  |A_R| \;\le\; n + \fprate(N - n) \;\approx\; n + \fprate N \quad (N \gg n).
\end{equation}
The verdict on $x$ depends on $R$ only through which side of $A_R$ the query
lands, so $A_R$ is the representation-independent object that does the work.

\subsection*{The theorem}

\begin{theorem}[Entropy floor for FNR-zero approximate sets]
\label{thm:entropy-floor}
Let $(R, \mathrm{query})$ be an FNR-zero representation scheme of bit length $L$
that, for every $n$-element $A \subseteq \Universe$, produces $A_R \supseteq A$
with $|A_R \setminus A| \le \fprate(N - n)$. Then in the regime $N \gg n$ with
$\fprate \in (0,1)$ bounded away from $1$,
\begin{equation}
\label{eq:entropy-floor}
  L \;\ge\; n \log_2(1/\fprate) - o(n),
  \qquad\text{equivalently}\qquad
  L/n \;\ge\; -\log_2 \fprate + o(1).
\end{equation}
\end{theorem}

\begin{proof}
The argument is a pigeonhole count. Fix the accepted set $A_R$ of a
representation $R$. By FNR-zero $A \subseteq A_R$, and by
\eqref{eq:accepted-size} $|A_R| \le n + \fprate N$, so the true set is some
$n$-subset of $A_R$: any of $\binom{|A_R|}{n}$ choices is consistent with what
$R$ reports. The algorithm cannot distinguish among them, so $R$ describes an
equivalence class of that size, not a single $A$. Every one of the
$\binom{N}{n}$ possible target sets needs some $R \in \{0,1\}^L$ whose accepted
set contains it, and there are at most $2^L$ strings, each covering at most
$\max_{A_R}\binom{|A_R|}{n}$ sets, so
\begin{equation}
\label{eq:pigeonhole}
  2^L \cdot \max_{A_R}\binom{|A_R|}{n} \;\ge\; \binom{N}{n}.
\end{equation}
Take logarithms, bound $|A_R| \le n + \fprate N$, and apply Stirling:
$\log_2\binom{N}{n} = n\log_2(N/n) + O(n)$ and, in the regime $\fprate N \gg
n$, $\log_2\binom{n + \fprate N}{n} = n\log_2(\fprate N/n) + O(n)$. Subtracting,
\[
  L \;\ge\; n\log_2(N/n) - n\log_2(\fprate N/n) - O(n)
       \;=\; n\log_2(1/\fprate) - O(n),
\]
and dividing by $n$ gives $L/n \ge -\log_2 \fprate + o(1)$. The full asymptotic
control of the $o(n)$ term is in \textcite{towell-bernoulli-entropy}.
\end{proof}

The simplification $\fprate N \gg n$ is the deployment regime: $\fprate$ fixed,
$N$ large relative to $n$, the target embedded in a much larger universe. There
the floor is sharp; when $\Universe$ is barely larger than $A$ the bound holds
only in the rougher form \eqref{eq:pigeonhole} and the per-element reading
degrades. (At the other extreme $\fprate = 1$ the floor is $0$, met by the
trivial always-accept structure that stores nothing.)

\subsection*{Form of the bound}

The bound has no parameters to tune: the per-element cost is a logarithm of the
inverse error rate, independent of how the representation is built.
\begin{center}
\begin{tabular}{l|ll}
\hline
  $\fprate$ & $-\log_2 \fprate$ (bits/element) & bits for $n = 10^6$ \\
\hline
  $10^{-1}$ & $\approx 3.32$  & $\approx 0.42$\,MB \\
  $10^{-2}$ & $\approx 6.64$  & $\approx 0.83$\,MB \\
  $10^{-3}$ & $\approx 9.97$  & $\approx 1.25$\,MB \\
  $10^{-6}$ & $\approx 19.93$ & $\approx 2.49$\,MB \\
  $10^{-9}$ & $\approx 29.90$ & $\approx 3.74$\,MB \\
\hline
\end{tabular}
\end{center}
A million-element FNR-zero set at one-in-a-billion false-positive rate cannot
fit in less than about $3.74$ megabytes, whatever the encoding. The cost grows
\emph{additively}, by $\log_2 10 \approx 3.32$ bits per element for each order
of magnitude in $1/\fprate$, not geometrically. The universe size $N$ does not
enter the leading term: one might have expected the $\log_2\binom{N}{n}$ cost
of identifying $A$ exactly, but the FPR budget buys back precisely that
universe-dependent excess, leaving the residual $n\log_2(1/\fprate)$.

\subsection*{Why the floor is one-sided}

\Cref{thm:entropy-floor} is asymmetric because its constraint is. The
pigeonhole step \eqref{eq:pigeonhole} uses FNR-zero to force $A \subseteq A_R$,
and that containment is what makes $\binom{|A_R|}{n}$ the right
equivalence-class count. If the false-negative rate were also positive, a true
element could be misclassified, $A$ would no longer sit inside $A_R$, and the
containment the proof rests on would break.

The two-sided lower bound is, at present, an open problem. A natural heuristic
keeps a $\tprate = 1 - \fnrate$ fraction of the $n$ elements on average, each
still costing $-\log_2 \fprate$ bits, which suggests a cost near $-(1 -
\fnrate)\,n\log_2 \fprate$; but this is \emph{not} a valid information-theoretic
lower bound. Randomly dropping elements changes the encoding problem
qualitatively rather than merely scaling it, and the pigeonhole count no longer
applies. What can be said is only qualitative: a two-sided set may in principle
be represented below the FNR-zero floor $n\log_2(1/\fprate)$, spending part of
its error budget on false negatives to blur the identity of $A$, but no proven
bound quantifies how far below. The FNR-zero floor is the rigorous result, and
it is the one the rest of the book builds on.

\subsection*{The constructive question}

\Cref{thm:entropy-floor} forbids: no FNR-zero scheme beats $-\log_2 \fprate$
bits per element. Whether the floor is achievable is the dual question. The
standard Bloom filter pays a factor of about $1.44$, derived in
\Cref{sec:bloom-gap}; chapter~14 (\Cref{ch:optimal-bernoulli-sets}) constructs
the Bernoulli Hash Function, which meets the floor asymptotically. The two
together characterize the FNR-zero family tightly: the floor sits at $-\log_2
\fprate$, and a known scheme reaches it. Until then the floor is a target, and
\Cref{sec:bloom-gap} measures how far the standard Bloom filter sits above it
and why.
